\chapter{Introduction}

\section{Ricci flow: what is it, and from where did it come?}

Let $\mathcal M$ be a smooth closed manifold equipped with a smooth Riemannian
metric $g$. Ricci flow evolves $g$ by the partial differential equation
of curvature form below.

\begin{definition}[Ricci flow evolution equation]
\label{topping-intro-ricci-flow-evolution}
\dcref{1.1.1}
\group{ricci-flow}
\source{topping:page-0007}
\uses{topping-intro-ricci-curvature}
The Ricci flow equation is
\[
\frac{\partial g}{\partial t} = -2\,\Ric(g).
\]
\end{definition}

\begin{remark}
\label{topping-intro-ricci-curvature}
\dcref{Ricci curvature}
\group{ricci-flow}
\source{topping:page-0007}
Here $\Ric(g)$ denotes the Ricci curvature of the metric $g$.
\end{remark}

In simple situations the flow deforms a metric into one distinguished by its
curvature. In dimension two, after a suitable renormalisation, it deforms a
metric conformally to one of constant curvature and so yields a proof of the
two-dimensional uniformisation theorem. More generally, topology may prevent
the existence of such distinguished metrics, and the flow can then develop a
finite-time singularity. The strategy is to stop at a singularity and perform
surgeries that remove the singular regions before continuing the flow.

Provided the singularity structure is sufficiently understood, one may hope to
reconstruct the topology of the original manifold from the limiting flow and
the pieces removed during surgery. These notes develop some of the machinery
used to analyse singularities and illustrate it by proving Hamilton's theorem
that closed three-manifolds with positive Ricci curvature admit a metric of
constant positive sectional curvature.

One motivation for the choice of evolution equation is that it behaves like a
heat equation for the metric. In harmonic coordinates the Ricci tensor has the
form
\[
R_{ij} = -\frac12 \Delta g_{ij} + \text{lower order terms},
\]
while in normal coordinates centred at a point $p$ one has
\[
R_{ij} = -\frac32 \Delta g_{ij}
\]
at $p$. Here the notation $\Delta g_{ij}$ is to be read carefully: it does not
mean the rough Laplacian applied to the metric tensor, which vanishes because
the metric is parallel.

\section{Examples and special solutions}

\subsection{Einstein manifolds}

A round sphere evolves by shrinking homothetically to a point in finite time.

\begin{remark}
\label{topping-intro-einstein-flow}
\dcref{Einstein metric flow}
\group{examples}
\source{topping:page-0009}
If $\Ric(g_0)=\lambda g_0$ for some constant $\lambda\in\mathbb R$, then
\[
g(t)=(1-2\lambda t)g_0
\]
is a solution of Ricci flow with initial condition $g(0)=g_0$.
\end{remark}

\begin{proof}
Differentiate $g(t)$ to get $\partial_t g(t)=-2\lambda g_0$. Since
\[
\Ric(g_0)=\lambda g_0,
\]
and the Ricci tensor is invariant under uniform scaling,
\[
-2\Ric(g(t))=-2\Ric((1-2\lambda t)g_0)=-2\Ric(g_0)=-2\lambda g_0,
\]
which matches $\partial_t g(t)$.
\end{proof}

In particular, for the round unit sphere $(S^n,g_0)$ one has
\[
\Ric(g_0)=(n-1)g_0,
\]
so the flow is $g(t)=(1-2(n-1)t)g_0$ and the sphere collapses at
\[
T=\frac{1}{2(n-1)}.
\]

If $g_0$ is hyperbolic, meaning that it has constant sectional curvature $-1$,
then $\Ric(g_0)=-(n-1)g_0$ and the flow expands homothetically for all time.

\subsection{Ricci solitons}

There is a more general notion of self-similarity than uniform shrinking or
expansion. Let $X(t)$ be a time-dependent smooth vector field on $\mathcal M$
with flow $\psi_t$. For any smooth $f:\mathcal M\to\mathbb R$,
\begin{definition}[Flow of a time-dependent vector field]
\label{topping-intro-flow-field}
\dcref{1.2.1}
\group{ricci-solitons}
\source{topping:page-0009}
The flow satisfies
\[
X(\psi_t(y),t)f = \frac{\partial}{\partial t}(f\circ\psi_t)(y).
\]
\end{definition}

\begin{proposition}
\label{topping-intro-pushforward-evolution}
\dcref{1.2.1--1.2.3}
\group{ricci-solitons}
\source{topping:page-0010}
\uses{topping-intro-flow-field}
If $\sigma(t)$ is a smooth function and
\[
\hat g(t)=\sigma(t)\psi_t^*(g(t)),
\]
then
\[
\frac{\partial \hat g}{\partial t}
=\sigma'(t)\psi_t^*(g)+\sigma(t)\psi_t^*\!\left(\frac{\partial g}{\partial t}\right)
+\sigma(t)\psi_t^*(\mathcal L_X g).
\]
\end{proposition}

\begin{proof}
Differentiate the product $\sigma(t)\psi_t^*(g(t))$ using the ordinary product
rule. The derivative of the pullback term is the Lie derivative term by the
definition of $\mathcal L_X$ and the characterisation in
\eqref{topping-intro-flow-field}.
\end{proof}

\begin{definition}
\label{topping-intro-ricci-soliton}
\dcref{1.2.2--1.2.4}
\group{ricci-solitons}
\source{topping:page-0010}
If a metric $g_0$, a vector field $Y$, and a constant $\lambda\in\mathbb R$
satisfy
\[
-2\Ric(g_0)=\mathcal L_Y g_0-2\lambda g_0,
\]
then the associated self-similar Ricci flow is called a steady, expanding, or
shrinking Ricci soliton according to whether $\lambda=0$, $\lambda<0$, or
$\lambda>0$.
\end{definition}

\begin{definition}
\label{topping-intro-gradient-soliton}
\dcref{1.2.3--1.2.5}
\group{ricci-solitons}
\source{topping:page-0011}
A Ricci soliton whose vector field is the gradient of a function
$f:\mathcal M\to\mathbb R$ is called a gradient Ricci soliton. In this case
\[
\Hess_{g_0}(f)+\Ric(g_0)=\lambda g_0.
\]
\end{definition}

\begin{example}
\label{topping-intro-cigar-soliton}
\dcref{Hamilton's cigar soliton}
\group{ricci-solitons}
\source{topping:page-0011}
On $\mathbb R^2$, let $g_0=\rho^2(dx^2+dy^2)$ with
\[
\rho^2=\frac{1}{1+x^2+y^2}.
\]
Then
\[
\Ric(g_0)=\frac{2}{1+x^2+y^2}g_0,
\]
and for
\[
Y=-2\!\left(x\frac{\partial}{\partial x}+y\frac{\partial}{\partial y}\right)
\]
one has
\[
\mathcal L_Y g_0=-\frac{4}{1+x^2+y^2}g_0.
\]
Hence $g_0$ is a steady Ricci soliton. In geodesic polar coordinates
\[
g_0=ds^2+\tanh^2(s)\,d\theta^2,
\]
with Gauss curvature
\[
K=\frac{2}{\cosh^2 s}.
\]
It is also a gradient soliton, with potential function
\[
f=-2\ln(\cosh s).
\]
\end{example}

The cigar is historically important because it appears as a possible blow-up
limit for three-dimensional Ricci flows near finite-time singularities. Later
work of Perelman rules out that possibility. Another rotationally symmetric
steady gradient soliton exists on $\mathbb R^3$, due to Bryant; it opens
asymptotically like a paraboloid and has positive sectional curvature. The flat
metric on $\mathbb R^n$ gives the Gaussian soliton and can be viewed as steady,
expanding, or shrinking according to the choice of potential.

\subsection{Parabolic rescaling of Ricci flows}

\begin{remark}
\label{topping-intro-parabolic-rescaling}
\dcref{1.2.7--1.2.9}
\group{rescaling}
\source{topping:page-0012}
If $g(t)$ is a Ricci flow on $[0,T]$ and $\lambda>0$, then
\[
\hat g(x,t)=\lambda\, g(x,t/\lambda)
\]
is also a Ricci flow on $[0,\lambda T]$. Under this scaling the Ricci tensor is
invariant, while sectional curvature and scalar curvature scale by
$\lambda^{-1}$.
\end{remark}

The connection is unchanged under this parabolic rescaling.

\section{Getting a feel for Ricci flow}

The main use of this rescaling is to analyse singularities by blowing up a flow
where curvature becomes large and passing to a limit. This is a standard
strategy in geometric analysis.

\subsection{Two dimensions}

In two dimensions one has $\Ric(g)=Kg$. The flow therefore expands regions with
negative Gauss curvature and shrinks regions with positive Gauss curvature.

\subsection{Three dimensions}

The three-dimensional picture is more complicated. A neck pinch is the basic
heuristic singularity: the flow tends to contract a positively curved
cross-sectional sphere while stretching in the neck direction. A blow-up of
such a singularity is expected to look like a cylinder $S^2\times\mathbb R$,
and in some asymmetrical situations one expects a degenerate neck pinch whose
blow-up resembles the Bryant soliton.

\section*{Infinite time behaviour}

Ricci flow need not converge as $t\to\infty$. Hyperbolic metrics expand
homothetically, and solitons may evolve in more complicated ways. In
long-time behaviour one expects thick-thin decompositions and, in appropriate
settings, the emergence of geometric pieces.

\section{The topology and geometry of manifolds in low dimensions}

Let us consider only closed, oriented manifolds in this section.

\subsection{Dimension 1}

There is only one such manifold: the circle $S^1$.

\subsection{Dimension 2}

Closed surfaces are classified by their genus. Each admits a conformally
equivalent metric of constant Gauss curvature, and after rescaling the
curvature can be taken to be $1$, $0$, or $-1$. The universal cover is then
$S^2$, $\mathbb R^2$, or $\mathbb H^2$ respectively, yielding the classical
uniformisation theorem for compact Riemann surfaces.

For later comparison with dimension three, note that $\pi_1(S^2)=1$,
\[
\pi_1(T^2)=\mathbb Z\oplus\mathbb Z,
\]
and for higher genus surfaces the fundamental group is infinite but contains
no copy of $\mathbb Z\oplus\mathbb Z$.

\subsection{Dimension 3}

Thurston conjectured a classification for three-manifolds analogous to the
two-dimensional picture. For irreducible closed oriented three-manifolds, one
expects either a spherical space form, the presence of
$\mathbb Z\oplus\mathbb Z$ in the fundamental group, or a hyperbolic quotient.
Prime decomposition reduces the general case to irreducible pieces. Thurston
then introduces geometric pieces built from one of eight geometries:
\[
S^3,\ \mathbb R^3,\ \mathbb H^3,\ S^2\times\mathbb R,\ \mathbb H^2\times\mathbb R,\ Nil,\ Sol,\ \overline{\mathrm{SL}_2}(\mathbb R).
\]

\begin{definition}
\label{topping-intro-geometry}
\dcref{1.4.1}
\group{three-manifolds}
\source{topping:page-0021}
A geometry is a simply connected homogeneous unimodular Riemannian manifold.
\end{definition}

Here homogeneous means that any two points are related by an isometry, and
unimodular means that the geometry admits a discrete group of isometries with
compact quotient.

\begin{definition}
\label{topping-intro-geometric-manifold}
\dcref{1.4.2}
\group{three-manifolds}
\source{topping:page-0021}
A compact manifold $\mathcal M$ is geometric if $\Int(\mathcal M)=X/\Gamma$
has finite volume, where $X$ is a geometry and $\Gamma$ is a discrete group of
isometries acting freely.
\end{definition}

When a geometric manifold has boundary, that boundary is a union of
incompressible tori.

\begin{conjecture}
\label{topping-intro-thurston-geometrization}
\dcref{1.4.3}
\group{three-manifolds}
\source{topping:page-0021}
Every smooth, closed, oriented prime three-manifold can be obtained by gluing a
finite number of geometric pieces along their boundary tori.
\end{conjecture}

Hamilton's theorem, proved later in these notes, gives a special case of this
conjecture for manifolds of positive Ricci curvature.

\section{Using Ricci flow to prove topological and geometric results}

In dimension two, the Ricci flow proves the uniformisation theorem after
renormalisation. In dimension three, the long-term program for the geometrisation
conjecture begins with arbitrary metrics, expects singularities such as neck
pinches, and uses surgery to excise the necks and restart the flow. In
favourable cases, such as positive Ricci curvature, the flow becomes extinct
and one can conclude that the manifold is a spherical space form. In general,
one expects a thick-thin decomposition at large times. Four-dimensional analogues
exist under additional curvature hypotheses, but are substantially harder.
